\documentclass[a4paper,13pt]{extreport}
\usepackage{extsizes}
\usepackage[utf8]{inputenc}
\usepackage[T5]{fontenc}
\usepackage{mathptmx}
\usepackage[vietnamese]{babel}
\usepackage{graphicx}
\usepackage{geometry}
\usepackage{titlesec}
\usepackage{fancyhdr}
\usepackage{hyperref}
\usepackage{listings}
\usepackage{xcolor}
\usepackage{float}
\usepackage{array}
\usepackage{longtable}
\usepackage{booktabs}
\usepackage{enumitem}
\usepackage{indentfirst}
\usepackage{tocloft}

\geometry{left=3cm,right=2cm,top=2cm,bottom=2cm}
\setlength{\parindent}{1.25cm} % Cấu hình thụt đầu dòng 1.25cm

% Đường dẫn ảnh - Overleaf sẽ tìm ảnh trong thư mục images/
\graphicspath{{images/}}

% Định dạng tiêu đề chương
\titleformat{\chapter}[display]
{\normalfont\Large\bfseries\centering}
{\MakeUppercase{\chaptertitlename}\ \thechapter}{0pt}{\Large\MakeUppercase}
\titlespacing*{\chapter}{0pt}{-20pt}{20pt}

% Header và Footer
\pagestyle{fancy}
\fancyhf{}
\fancyhead[R]{\thepage}
\renewcommand{\headrulewidth}{0pt}

% Định dạng code
\lstset{
    basicstyle=\ttfamily\small,
    breaklines=true,
    frame=single,
    backgroundcolor=\color{gray!10}
}

\begin{document}

% ==================== TRANG BÌA ====================
\begin{titlepage}
\begin{center}
% Logo trường
\includegraphics[width=3cm]{logo_tvu}\\[0.3cm]
\textbf{TRƯỜNG ĐẠI HỌC TRÀ VINH}\\
\textbf{KHOA KỸ THUẬT VÀ CÔNG NGHỆ}\\[0.5cm]
\rule{\linewidth}{0.5mm}\\[1cm]

\textbf{\Large ĐỒ ÁN CHUYÊN NGÀNH}\\[1.5cm]

\textbf{\LARGE XÂY DỰNG WEBSITE KẾT NỐI Y TẾ}\\
\textbf{\LARGE GIỮA BỆNH NHÂN VÀ SINH VIÊN Y KHOA}\\[2cm]

\begin{flushleft}
\textbf{Sinh viên thực hiện:} Trầm Tấn Khá\\[0.3cm]
\textbf{Mã số sinh viên:} 110122087\\[0.3cm]
\textbf{Giảng viên hướng dẫn:} Nguyễn Khắc Quốc\\[0.3cm]
\end{flushleft}

\vfill
\textbf{Trà Vinh, năm 2025}
\end{center}
\end{titlepage}

% ==================== NHẬN XÉT CỦA GVHD ====================
\newpage
\thispagestyle{empty}
\begin{center}
\textbf{\Large NHẬN XÉT CỦA GIẢNG VIÊN HƯỚNG DẪN}
\end{center}
\vspace{1cm}
\noindent\rule{\linewidth}{0.4pt}\\[0.5cm]
\noindent\rule{\linewidth}{0.4pt}\\[0.5cm]
\noindent\rule{\linewidth}{0.4pt}\\[0.5cm]
\noindent\rule{\linewidth}{0.4pt}\\[0.5cm]
\noindent\rule{\linewidth}{0.4pt}\\[0.5cm]
\noindent\rule{\linewidth}{0.4pt}\\[0.5cm]
\noindent\rule{\linewidth}{0.4pt}\\[0.5cm]
\noindent\rule{\linewidth}{0.4pt}\\[0.5cm]
\noindent\rule{\linewidth}{0.4pt}\\[0.5cm]
\noindent\rule{\linewidth}{0.4pt}\\[0.5cm]
\noindent\rule{\linewidth}{0.4pt}\\[0.5cm]
\noindent\rule{\linewidth}{0.4pt}\\[0.5cm]

\vspace{2cm}
\begin{flushright}
Trà Vinh, ngày \ldots\ldots tháng \ldots\ldots năm 2025\\[1cm]
\textbf{Giảng viên hướng dẫn}\\
(Ký tên và ghi rõ họ tên)\\[2cm]
\textbf{Nguyễn Khắc Quốc}
\end{flushright}

% ==================== NHẬN XÉT CỦA HỘI ĐỒNG ====================
\newpage
\thispagestyle{empty}
\begin{center}
\textbf{\Large NHẬN XÉT CỦA THÀNH VIÊN HỘI ĐỒNG}
\end{center}
\vspace{1cm}
\noindent\rule{\linewidth}{0.4pt}\\[0.5cm]
\noindent\rule{\linewidth}{0.4pt}\\[0.5cm]
\noindent\rule{\linewidth}{0.4pt}\\[0.5cm]
\noindent\rule{\linewidth}{0.4pt}\\[0.5cm]
\noindent\rule{\linewidth}{0.4pt}\\[0.5cm]
\noindent\rule{\linewidth}{0.4pt}\\[0.5cm]
\noindent\rule{\linewidth}{0.4pt}\\[0.5cm]
\noindent\rule{\linewidth}{0.4pt}\\[0.5cm]
\noindent\rule{\linewidth}{0.4pt}\\[0.5cm]
\noindent\rule{\linewidth}{0.4pt}\\[0.5cm]
\noindent\rule{\linewidth}{0.4pt}\\[0.5cm]
\noindent\rule{\linewidth}{0.4pt}\\[0.5cm]

\vspace{2cm}
\begin{flushright}
Trà Vinh, ngày \ldots\ldots tháng \ldots\ldots năm 2025\\[1cm]
\textbf{Thành viên hội đồng}\\
(Ký tên và ghi rõ họ tên)\\[2cm]
\end{flushright}

% ==================== LỜI CẢM ƠN ====================
\newpage
\thispagestyle{empty}
\begin{center}
\textbf{\Large LỜI CẢM ƠN}
\end{center}
\vspace{1cm}

Để hoàn thành đồ án chuyên ngành này, em xin gửi lời cảm ơn chân thành đến:

Thầy \textbf{Nguyễn Khắc Quốc} - Giảng viên hướng dẫn đã tận tình chỉ bảo, định hướng và hỗ trợ em trong suốt quá trình thực hiện đồ án.

Quý Thầy Cô trong Khoa Kỹ thuật và Công nghệ, Trường Đại học Trà Vinh đã truyền đạt những kiến thức quý báu trong suốt thời gian em học tập tại trường.

Gia đình và bạn bè đã luôn động viên, hỗ trợ em trong quá trình học tập và thực hiện đồ án.

Mặc dù đã cố gắng hoàn thiện đồ án, nhưng do kiến thức và kinh nghiệm còn hạn chế nên không tránh khỏi những thiếu sót. Em rất mong nhận được sự góp ý của Quý Thầy Cô để đồ án được hoàn thiện hơn.

Em xin chân thành cảm ơn!

\begin{flushright}
\textit{Trà Vinh, tháng 11 năm 2025}\\
\textbf{Sinh viên thực hiện}\\[1cm]
\textbf{Trầm Tấn Khá}
\end{flushright}

% ==================== MỤC LỤC ====================
\newpage
\tableofcontents

% ==================== DANH MỤC HÌNH ẢNH ====================
\newpage
\listoffigures
\addcontentsline{toc}{chapter}{DANH MỤC HÌNH ẢNH}

% ==================== DANH MỤC BẢNG BIỂU ====================
\newpage
\listoftables
\addcontentsline{toc}{chapter}{DANH MỤC BẢNG BIỂU}


% ==================== TÓM TẮT ĐỒ ÁN ====================
\newpage
\addcontentsline{toc}{chapter}{TÓM TẮT ĐỒ ÁN}
\begin{center}
\textbf{\Large TÓM TẮT ĐỒ ÁN CHUYÊN NGÀNH}
\end{center}
\vspace{1cm}

Đồ án "Xây dựng Website kết nối y tế giữa bệnh nhân và sinh viên Y khoa" được thực hiện nhằm giải quyết vấn đề kết nối giữa bệnh nhân có nhu cầu chăm sóc sức khỏe tại nhà với sinh viên Y khoa đang cần thực hành và tích lũy kinh nghiệm.

\textbf{Vấn đề nghiên cứu:} Hiện nay, nhiều bệnh nhân cần được chăm sóc y tế tại nhà nhưng gặp khó khăn trong việc tìm kiếm người hỗ trợ phù hợp. Đồng thời, sinh viên Y khoa cần cơ hội thực hành để nâng cao kỹ năng chuyên môn.

\textbf{Hướng tiếp cận:} Xây dựng một nền tảng web cho phép:
\begin{itemize}
    \item Bệnh nhân đăng tin tuyển dụng người chăm sóc
    \item Sinh viên Y khoa đăng tin ứng tuyển
    \item Hai bên có thể trao đổi, liên lạc trực tiếp qua hệ thống tin nhắn
    \item Hệ thống xác thực sinh viên đảm bảo độ tin cậy
\end{itemize}

\textbf{Giải pháp kỹ thuật:}
\begin{itemize}
    \item Ngôn ngữ lập trình: PHP
    \item Cơ sở dữ liệu: MySQL
    \item Giao diện: HTML, CSS, Bootstrap
    \item Môi trường phát triển: XAMPP
\end{itemize}

\textbf{Kết quả đạt được:}
\begin{itemize}
    \item Hệ thống đăng ký, đăng nhập với phân quyền người dùng
    \item Chức năng đăng tin tuyển dụng và ứng tuyển
    \item Hệ thống tin nhắn nội bộ
    \item Trang quản trị với đầy đủ chức năng quản lý
    \item Xác thực sinh viên qua thẻ sinh viên
    \item Hệ thống đánh giá và báo cáo vi phạm
\end{itemize}

% ==================== MỞ ĐẦU ====================
\newpage
\addcontentsline{toc}{chapter}{MỞ ĐẦU}
\begin{center}
\textbf{\Large MỞ ĐẦU}
\end{center}
\vspace{1cm}

\section*{1. Lý do chọn đề tài}
\addcontentsline{toc}{section}{1. Lý do chọn đề tài}

Trong bối cảnh xã hội ngày càng phát triển cùng với tốc độ già hóa dân số nhanh chóng, nhu cầu chăm sóc sức khỏe tại nhà ngày càng tăng cao, đặc biệt là đối với người cao tuổi, người mắc bệnh mãn tính hoặc bệnh nhân đang trong giai đoạn phục hồi chức năng sau phẫu thuật. Mặc dù vậy, việc tìm kiếm các dịch vụ hỗ trợ y tế tại gia chất lượng cao, có chuyên môn nghiệp vụ nhưng chi phí hợp lý vẫn là một bài toán khó đối với nhiều gia đình. Ở chiều ngược lại, sinh viên y khoa tại các trường đại học y dược luôn có nhu cầu tìm kiếm cơ hội thực hành lâm sàng ngoài giờ học. Việc chăm sóc bệnh nhân thực tế ngoài môi trường bệnh viện không chỉ giúp sinh viên tích lũy kinh nghiệm thực tiễn, hoàn thiện các kỹ năng chuyên môn nghiệp vụ và kỹ năng giao tiếp y học, mà còn mang lại nguồn thu nhập trang trải cuộc sống trong quá trình học tập.

Mặc dù nhu cầu kết nối giữa bệnh nhân và sinh viên y khoa là rất lớn và có tính tương hỗ cao, song phương thức kết nối hiện nay chủ yếu diễn ra tự phát trên các nền tảng mạng xã hội không chuyên dụng (như Facebook, Zalo, các hội nhóm cộng đồng). Điều này dẫn đến nhiều hạn chế nghiêm trọng: khó kiểm chứng danh tính và năng lực chuyên môn của sinh viên, khó sàng lọc đối tượng phù hợp với bệnh lý cụ thể, và thiếu tính minh bạch, an toàn về mặt thông tin cũng như pháp lý cho cả hai bên. Bên cạnh đó, đối với các nhu cầu y tế cơ bản (như đo huyết áp, tiêm truyền theo chỉ định, thay băng, rửa vết thương, tư vấn chăm sóc sức khỏe ban đầu), việc lựa chọn một sinh viên y khoa hỗ trợ tại nhà sẽ giúp bệnh nhân giảm thiểu thời gian chờ đợi tại bệnh viện, tiết kiệm chi phí đi lại và giảm tải đáng kể cho hệ thống y tế công cộng.

Do đó, đề tài ``Xây dựng Website kết nối y tế giữa bệnh nhân và sinh viên Y khoa'' được thực hiện nhằm xây dựng một nền tảng trực tuyến chuyên biệt, an toàn và hiệu quả. Ứng dụng cung cấp giải pháp kết nối hai chiều minh bạch nhờ quy trình xác minh thẻ sinh viên nghiêm ngặt kết hợp kiểm duyệt của quản trị viên, hỗ trợ giao tiếp đa kênh trực quan (nhắn tin nội bộ, gọi Video Call thời gian thực qua WebRTC) cùng trợ lý ảo AI hỗ trợ tư vấn y học cơ bản 24/7. Đây không chỉ là giải pháp tối ưu giúp bệnh nhân tiếp cận dịch vụ chăm sóc sức khỏe uy tín tại nhà với chi phí hợp lý, mà còn tạo môi trường thực hành lâm sàng thực tế cho sinh viên y khoa, hướng tới xây dựng một cộng đồng y tế tương hỗ, phát triển bền vững.

\section*{2. Mục đích nghiên cứu}
\addcontentsline{toc}{section}{2. Mục đích nghiên cứu}

\begin{itemize}
    \item Xây dựng website kết nối trực tuyến hai chiều giữa bệnh nhân có nhu cầu hỗ trợ y tế tại nhà với sinh viên Y khoa đang cần thực hành lâm sàng thực tế.
    \item Thiết lập cơ chế kiểm duyệt thẻ sinh viên và đánh giá hai chiều nhằm tạo dựng một không gian giao dịch y tế an toàn, đáng tin cậy.
    \item Tích hợp các giải pháp giao tiếp đa kênh bao gồm nhắn tin trực tuyến và gọi cuộc gọi Video thời gian thực (WebRTC Video Call).
    \item Nghiên cứu tích hợp Trợ lý ảo AI Chatbot hỗ trợ tư vấn các kiến thức y khoa cơ bản và định hướng dịch vụ chăm sóc sức khỏe.
    \item Tự động hóa quy trình quản lý lịch hẹn y tế (Appointments) giúp tối ưu thời gian và nâng cao trải nghiệm người dùng.
\end{itemize}

\section*{3. Đối tượng và phạm vi nghiên cứu}
\addcontentsline{toc}{section}{3. Đối tượng và phạm vi nghiên cứu}

\textbf{Đối tượng nghiên cứu:}
\begin{itemize}
    \item Bệnh nhân có nhu cầu tìm người chăm sóc y tế tại nhà
    \item Sinh viên Y khoa muốn tìm cơ hội thực hành và làm thêm
    \item Các công nghệ web: PHP, MySQL, HTML/CSS, Bootstrap
\end{itemize}

\textbf{Phạm vi nghiên cứu:}
\begin{itemize}
    \item Xây dựng website với các chức năng cơ bản: đăng ký, đăng nhập, đăng tin, tìm kiếm, nhắn tin
    \item Hệ thống quản trị cơ bản
    \item Triển khai trên môi trường localhost (XAMPP)
\end{itemize}

% ==================== CHƯƠNG 1: TỔNG QUAN ====================
\chapter{TỔNG QUAN}

\section{Lý do chọn đề tài}
Trong bối cảnh xã hội ngày càng phát triển cùng với tốc độ già hóa dân số nhanh chóng, nhu cầu chăm sóc sức khỏe tại nhà ngày càng tăng cao, đặc biệt là đối với người cao tuổi, người mắc bệnh mãn tính hoặc bệnh nhân đang trong giai đoạn phục hồi chức năng sau phẫu thuật. Mặc dù vậy, việc tìm kiếm các dịch vụ hỗ trợ y tế tại gia chất lượng cao, có chuyên môn nghiệp vụ nhưng chi phí hợp lý vẫn là một bài toán khó đối với nhiều gia đình. Ở chiều ngược lại, sinh viên y khoa tại các trường đại học y dược luôn có nhu cầu tìm kiếm cơ hội thực hành lâm sàng ngoài giờ học. Việc chăm sóc bệnh nhân thực tế ngoài môi trường bệnh viện không chỉ giúp sinh viên tích lũy kinh nghiệm thực tiễn, hoàn thiện các kỹ năng chuyên môn nghiệp vụ và kỹ năng giao tiếp y học, mà còn mang lại nguồn thu nhập trang trải cuộc sống trong quá trình học tập.

Mặc dù nhu cầu kết nối giữa bệnh nhân và sinh viên y khoa là rất lớn và có tính tương hỗ cao, song phương thức kết nối hiện nay chủ yếu diễn ra tự phát trên các nền tảng mạng xã hội không chuyên dụng (như Facebook, Zalo, các hội nhóm cộng đồng). Điều này dẫn đến nhiều hạn chế nghiêm trọng: khó kiểm chứng danh tính và năng lực chuyên môn của sinh viên, khó sàng lọc đối tượng phù hợp với bệnh lý cụ thể, và thiếu tính minh bạch, an toàn về mặt thông tin cũng như pháp lý cho cả hai bên. Bên cạnh đó, đối với các nhu cầu y tế cơ bản (như đo huyết áp, tiêm truyền theo chỉ định, thay băng, rửa vết thương, tư vấn chăm sóc sức khỏe ban đầu), việc lựa chọn một sinh viên y khoa hỗ trợ tại nhà sẽ giúp bệnh nhân giảm thiểu thời gian chờ đợi tại bệnh viện, tiết kiệm chi phí đi lại và giảm tải đáng kể cho hệ thống y tế công cộng.

Do đó, đề tài ``Xây dựng Website kết nối y tế giữa bệnh nhân và sinh viên Y khoa'' được thực hiện nhằm xây dựng một nền tảng trực tuyến chuyên biệt, an toàn và hiệu quả. Ứng dụng cung cấp giải pháp kết nối hai chiều minh bạch nhờ quy trình xác minh thẻ sinh viên nghiêm ngặt kết hợp kiểm duyệt của quản trị viên, hỗ trợ giao tiếp đa kênh trực quan (nhắn tin nội bộ, gọi Video Call thời gian thực qua WebRTC) cùng trợ lý ảo AI hỗ trợ tư vấn y học cơ bản 24/7. Đây không chỉ là giải pháp tối ưu giúp bệnh nhân tiếp cận dịch vụ chăm sóc sức khỏe uy tín tại nhà với chi phí hợp lý, mà còn tạo môi trường thực hành lâm sàng thực tế cho sinh viên y khoa, hướng tới xây dựng một cộng đồng y tế tương hỗ, phát triển bền vững.

\section{Mục đích nghiên cứu}
Mục đích chính của đề tài là xây dựng một website hoàn chỉnh hỗ trợ kết nối hai chiều trực tuyến giữa bệnh nhân (hoặc người thân có nhu cầu tìm kiếm người hỗ trợ y tế tại gia) và sinh viên Y khoa đang học tập tại các trường đại học y dược (cần cơ hội thực hành lâm sàng thực tế ngoài giờ học). Để giải quyết triệt để các rủi ro của phương thức kết nối tự phát truyền thống và nâng cao trải nghiệm người dùng, đề tài tập trung hiện thực hóa các mục tiêu cốt lõi sau:
\begin{itemize}
    \item \textbf{Xây dựng môi trường kết nối trực tuyến an toàn và đáng tin cậy:} Triển khai cơ chế xác minh tài khoản sinh viên qua thẻ sinh viên kết hợp quy trình duyệt hồ sơ thủ công của quản trị viên (Admin), từ đó kiểm chứng danh tính và ngăn chặn các tài khoản mạo danh, giả danh sinh viên Y khoa.
    \item \textbf{Tích hợp công cụ giao tiếp đa kênh trực quan thời gian thực:} Hiện thực hóa tính năng nhắn tin nội bộ trao đổi thông tin bảo mật và tích hợp giải pháp cuộc gọi Video trực tuyến (Video Call sử dụng công nghệ WebRTC). Tính năng này giúp bệnh nhân đối chiếu trực quan sinh viên Y khoa trước khi gặp mặt, đồng thời giúp sinh viên xem trước tình trạng lâm sàng (vết thương, đơn thuốc, bệnh án) của bệnh nhân để chuẩn bị dụng cụ y tế phù hợp.
    \item \textbf{Tích hợp Trợ lý ảo trí tuệ nhân tạo (AI Assistant) hỗ trợ 24/7:} Nghiên cứu và triển khai chatbot thông minh (kết nối qua Hugging Face/Gemini APIs) hỗ trợ tư vấn các kiến thức y khoa cơ bản, hướng dẫn sơ cứu ban đầu và giải đáp thắc mắc về quy trình vận hành hệ thống cho người dùng mọi lúc mọi nơi.
    \item \textbf{Thiết lập quy trình quản lý chuyên nghiệp và nâng cao uy tín cộng đồng:} Tự động hóa quy trình đặt lịch hẹn y tế (Appointments) ngay khi bệnh nhân phê duyệt sinh viên, kết hợp cơ chế đánh giá hai chiều (Rating \& Review) để duy trì chất lượng dịch vụ và nâng cao trách nhiệm của cộng đồng.
    \item \textbf{Hiện thực hóa về mặt kỹ thuật và nghiên cứu học thuật:} Ứng dụng các phương pháp phân tích thiết kế hệ thống chuẩn hóa (mô hình UML: Use Case Diagram, Class Diagram, ERD), xây dựng kiến trúc Client-Server dựa trên backend PHP thuần, hệ quản trị cơ sở dữ liệu MySQL, frontend responsive bằng Bootstrap 5.3, và ảo hóa container qua Docker nhằm tối ưu hóa tính mở rộng, bảo mật và khả năng triển khai thực tiễn của phần mềm.
\end{itemize}

\section{Đối tượng nghiên cứu}
Đối tượng nghiên cứu của đề tài tập trung vào các chủ thể và khía cạnh công nghệ sau:
\begin{enumerate}
    \item \textbf{Đối tượng người dùng}:
    \begin{itemize}
        \item Bệnh nhân hoặc người nhà bệnh nhân có nhu cầu tìm người chăm sóc sức khỏe tại nhà (đặc biệt là người cao tuổi, người mắc bệnh mãn tính hoặc người cần phục hồi chức năng sau phẫu thuật).
        \item Sinh viên Y khoa tại các trường đại học y dược có nhu cầu tìm kiếm công việc chăm sóc sức khỏe tại gia ngoài giờ học để tích lũy kinh nghiệm lâm sàng và có thêm thu nhập.
    \end{itemize}
    \item \textbf{Đối tượng công nghệ}:
    \begin{itemize}
        \item Quy trình và nghiệp vụ kết nối y tế (quy trình đăng tin tuyển dụng/ứng tuyển, quy trình duyệt ứng viên và tự động tạo lịch hẹn).
        \item Các công nghệ lập trình web: ngôn ngữ backend PHP, hệ quản trị cơ sở dữ liệu MySQL, frontend Bootstrap 5.3.
        \item Các công nghệ kết nối nâng cao: WebRTC cho cuộc gọi video Peer-to-Peer, tích hợp APIs trí tuệ nhân tạo (Gemini AI/Hugging Face).
        \item Môi trường phát triển và triển khai: XAMPP, Docker containerization.
    \end{itemize}
\end{enumerate}

\section{Phạm vi nghiên cứu}
Phạm vi nghiên cứu của đề tài được phân định như sau:
\begin{enumerate}
    \item \textbf{Phạm vi nội dung và chức năng}:
    \begin{itemize}
        \item Thiết kế và cài đặt các chức năng cốt lõi cho hai đối tượng người dùng chính: Bệnh nhân (đăng tin tuyển dụng kèm ảnh minh chứng bệnh án, tìm kiếm sinh viên, nhắn tin, gọi video, tạo lịch hẹn, đánh giá) và Sinh viên (đăng tin ứng tuyển kèm thẻ sinh viên, xin xác thực, kết bạn, nhắn tin, xem lịch sử nhận việc).
        \item Thiết kế trang quản trị (Admin Dashboard) hỗ trợ kiểm duyệt bài đăng, duyệt xác minh thẻ sinh viên, xử lý báo cáo vi phạm và gửi thông báo hệ thống.
        \item Tích hợp Chatbot AI hỗ trợ tư vấn y học cơ bản.
        \item \textit{Giới hạn}: Đồ án chưa tích hợp hệ thống thanh toán trực tuyến (e-payment) hoặc đặt cọc giữ tiền (escrow payment); giao dịch tài chính và thanh toán hiện tại do hai bên tự thỏa thuận ngoài hệ thống. Chưa xây dựng ứng dụng di động native (iOS/Android).
    \end{itemize}
    \item \textbf{Phạm vi kỹ thuật và môi trường}:
    \begin{itemize}
        \item Hệ thống được hiện thực hóa dưới dạng Website Responsive hoạt động tối ưu trên cả desktop và thiết bị di động.
        \item Môi trường phát triển và kiểm thử được triển khai cục bộ (Localhost) bằng XAMPP và ảo hóa container qua Docker (Apache + PHP + MySQL). Chưa triển khai trên các hệ thống Cloud server thực tế.
    \end{itemize}
\end{enumerate}


% ==================== CHƯƠNG 2: NGHIÊN CỨU LÝ THUYẾT ====================
\chapter{NGHIÊN CỨU LÝ THUYẾT}

\section{Tổng quan về phát triển Web}

\subsection{Khái niệm Website}
Website là tập hợp các trang web có liên quan với nhau, được lưu trữ trên máy chủ web và có thể truy cập thông qua Internet bằng trình duyệt web. Mỗi website có một địa chỉ URL duy nhất.

\subsection{Kiến trúc Client-Server}
Mô hình Client-Server là kiến trúc phổ biến trong phát triển web:
\begin{itemize}
    \item \textbf{Client (Máy khách):} Trình duyệt web của người dùng, gửi yêu cầu đến server
    \item \textbf{Server (Máy chủ):} Xử lý yêu cầu, truy vấn cơ sở dữ liệu và trả về kết quả
\end{itemize}

\section{Ngôn ngữ lập trình PHP}

\subsection{Giới thiệu PHP}
PHP (Hypertext Preprocessor) là ngôn ngữ lập trình kịch bản phía server, được sử dụng rộng rãi trong phát triển web. PHP có các đặc điểm:
\begin{itemize}
    \item Mã nguồn mở, miễn phí
    \item Dễ học, dễ sử dụng
    \item Hỗ trợ nhiều hệ quản trị cơ sở dữ liệu
    \item Cộng đồng lớn, nhiều tài liệu hỗ trợ
    \item Tương thích với hầu hết các web server
\end{itemize}

\subsection{Cú pháp cơ bản PHP}
PHP được nhúng trong HTML với cặp thẻ \texttt{<?php ?>}. Các đặc điểm cú pháp:
\begin{itemize}
    \item Biến bắt đầu bằng ký tự \$
    \item Hỗ trợ các kiểu dữ liệu: string, integer, float, boolean, array, object
    \item Cấu trúc điều khiển: if-else, switch, for, while, foreach
    \item Hỗ trợ lập trình hướng đối tượng
\end{itemize}

\section{Hệ quản trị cơ sở dữ liệu MySQL}

\subsection{Giới thiệu MySQL}
MySQL là hệ quản trị cơ sở dữ liệu quan hệ (RDBMS) mã nguồn mở phổ biến nhất. Đặc điểm:
\begin{itemize}
    \item Hiệu suất cao, ổn định
    \item Hỗ trợ nhiều nền tảng
    \item Bảo mật tốt
    \item Dễ dàng tích hợp với PHP
    \item Hỗ trợ giao dịch (transaction)
\end{itemize}

\subsection{Ngôn ngữ SQL}
SQL (Structured Query Language) là ngôn ngữ truy vấn cơ sở dữ liệu:
\begin{itemize}
    \item \textbf{DDL (Data Definition Language):} CREATE, ALTER, DROP
    \item \textbf{DML (Data Manipulation Language):} SELECT, INSERT, UPDATE, DELETE
    \item \textbf{DCL (Data Control Language):} GRANT, REVOKE
\end{itemize}

\section{Công nghệ Front-end}

\subsection{HTML (HyperText Markup Language)}
HTML là ngôn ngữ đánh dấu siêu văn bản, dùng để tạo cấu trúc trang web. HTML5 là phiên bản mới nhất với nhiều tính năng:
\begin{itemize}
    \item Các thẻ ngữ nghĩa: header, nav, section, article, footer
    \item Hỗ trợ multimedia: audio, video
    \item Form controls mới: date, email, number, range
    \item Canvas và SVG cho đồ họa
\end{itemize}

\subsection{CSS (Cascading Style Sheets)}
CSS dùng để định dạng giao diện trang web:
\begin{itemize}
    \item Tách biệt nội dung và trình bày
    \item Hỗ trợ responsive design
    \item CSS3 với animation, transition, flexbox, grid
\end{itemize}

\subsection{Bootstrap Framework}
Bootstrap là framework CSS phổ biến:
\begin{itemize}
    \item Hệ thống grid responsive
    \item Các component UI sẵn có
    \item Tương thích đa trình duyệt
    \item Dễ tùy chỉnh
\end{itemize}

\section{Môi trường phát triển XAMPP}

XAMPP là gói phần mềm tích hợp bao gồm:
\begin{itemize}
    \item \textbf{X:} Cross-platform (đa nền tảng)
    \item \textbf{A:} Apache web server
    \item \textbf{M:} MySQL/MariaDB database
    \item \textbf{P:} PHP
    \item \textbf{P:} Perl
\end{itemize}

XAMPP cung cấp môi trường phát triển web hoàn chỉnh trên máy local, phù hợp cho việc học tập và phát triển ứng dụng.

\section{Phương pháp phân tích thiết kế hệ thống}

\subsection{UML (Unified Modeling Language)}
UML là ngôn ngữ mô hình hóa thống nhất, sử dụng các loại biểu đồ:
\begin{itemize}
    \item Use Case Diagram: Mô tả chức năng hệ thống
    \item Class Diagram: Mô tả cấu trúc lớp
    \item Sequence Diagram: Mô tả tương tác theo thời gian
    \item Activity Diagram: Mô tả luồng hoạt động
    \item ERD (Entity Relationship Diagram): Mô tả cơ sở dữ liệu
\end{itemize}


% ==================== CHƯƠNG 3: HIỆN THỰC HÓA NGHIÊN CỨU ====================
\chapter{HIỆN THỰC HÓA NGHIÊN CỨU}

\section{Phân tích yêu cầu hệ thống}

\subsection{Yêu cầu chức năng}

\subsubsection{Đối với Bệnh nhân}
\begin{itemize}
    \item Đăng ký tài khoản với vai trò bệnh nhân
    \item Đăng nhập/Đăng xuất
    \item Xem danh sách tin ứng tuyển của sinh viên
    \item Đăng tin tuyển dụng người chăm sóc
    \item Tìm kiếm, lọc tin theo khu vực, danh mục
    \item Lưu tin yêu thích
    \item Gửi tin nhắn cho sinh viên
    \item Đánh giá sinh viên sau khi sử dụng dịch vụ
    \item Quản lý hồ sơ cá nhân
\end{itemize}

\subsubsection{Đối với Sinh viên Y khoa}
\begin{itemize}
    \item Đăng ký tài khoản với vai trò sinh viên
    \item Xác thực tài khoản bằng thẻ sinh viên
    \item Xem danh sách tin tuyển dụng
    \item Đăng tin ứng tuyển (sau khi được xác thực)
    \item Xin cấp quyền đăng tin
    \item Gửi tin nhắn cho bệnh nhân
    \item Quản lý hồ sơ cá nhân
    \item Xem lịch sử nhận việc
\end{itemize}

\subsubsection{Đối với Quản trị viên}
\begin{itemize}
    \item Quản lý người dùng (xem, khóa, xóa)
    \item Quản lý bài viết (duyệt, ẩn, xóa)
    \item Xử lý yêu cầu xác thực sinh viên
    \item Xử lý yêu cầu cấp quyền đăng tin
    \item Xem và xử lý báo cáo vi phạm
    \item Gửi thông báo đến người dùng
    \item Xem thống kê hệ thống
\end{itemize}

\subsection{Yêu cầu phi chức năng}
\begin{itemize}
    \item Giao diện thân thiện, dễ sử dụng
    \item Tương thích với nhiều trình duyệt
    \item Responsive trên các thiết bị
    \item Bảo mật thông tin người dùng
    \item Thời gian phản hồi nhanh
\end{itemize}

\section{Thiết kế hệ thống}

\subsection{Biểu đồ Use Case}

Biểu đồ Use Case mô tả các chức năng chính của hệ thống và các tác nhân tương tác.

% LƯU Ý: Cần chuyển file usecase.SVG sang usecase.png và upload vào thư mục images/
% Sử dụng công cụ online: https://svgtopng.com/ hoặc Inkscape
\begin{figure}[H]
    \centering
    \includegraphics[width=0.9\textwidth]{usecase.png}
    \caption{Biểu đồ Use Case tổng quan hệ thống}
    \label{fig:usecase}
\end{figure}

\textbf{Các tác nhân (Actors):}
\begin{itemize}
    \item \textbf{Khách (Guest):} Người dùng chưa đăng nhập
    \item \textbf{Bệnh nhân (Patient):} Người dùng có nhu cầu tìm người chăm sóc
    \item \textbf{Sinh viên (Student):} Sinh viên Y khoa muốn tìm việc
    \item \textbf{Quản trị viên (Admin):} Người quản lý hệ thống
\end{itemize}

\subsection{Thiết kế cơ sở dữ liệu}

\subsubsection{Biểu đồ ERD}

% LƯU Ý: Cần chuyển file erd.SVG sang erd.png và upload vào thư mục images/
% Sử dụng công cụ online: https://svgtopng.com/ hoặc Inkscape
\begin{figure}[H]
    \centering
    \includegraphics[width=0.95\textwidth]{erd.png}
    \caption{Biểu đồ ERD - Quan hệ thực thể}
    \label{fig:erd}
\end{figure}

\subsubsection{Mô tả các bảng dữ liệu}

\begin{table}[H]
\centering
\caption{Bảng Users - Thông tin người dùng}
\begin{tabular}{|l|l|l|}
\hline
\textbf{Tên trường} & \textbf{Kiểu dữ liệu} & \textbf{Mô tả} \\
\hline
id & INT (PK) & Mã người dùng \\
name & VARCHAR(150) & Họ tên \\
email & VARCHAR(255) & Email (unique) \\
email\_verified & TINYINT & Trạng thái xác thực email \\
password & VARCHAR(255) & Mật khẩu (hash) \\
role & ENUM & Vai trò (patient/student) \\
bio & TEXT & Giới thiệu bản thân \\
location & VARCHAR(255) & Địa chỉ \\
phone & VARCHAR(50) & Số điện thoại \\
student\_id & VARCHAR(100) & Mã số sinh viên \\
verified & TINYINT & Trạng thái xác thực SV \\
is\_admin & TINYINT & Quyền admin \\
can\_post & TINYINT & Quyền đăng tin \\
created\_at & TIMESTAMP & Ngày tạo \\
\hline
\end{tabular}
\end{table}

\begin{table}[H]
\centering
\caption{Bảng Posts - Bài đăng}
\begin{tabular}{|l|l|l|}
\hline
\textbf{Tên trường} & \textbf{Kiểu dữ liệu} & \textbf{Mô tả} \\
\hline
id & INT (PK) & Mã bài đăng \\
user\_id & INT (FK) & Mã người đăng \\
title & VARCHAR(255) & Tiêu đề \\
content & TEXT & Nội dung \\
type & ENUM & Loại (recruitment/application) \\
category & VARCHAR(100) & Danh mục \\
area & VARCHAR(255) & Khu vực \\
contact\_info & VARCHAR(255) & Thông tin liên hệ \\
suggested\_price & INT & Giá đề xuất \\
status & ENUM & Trạng thái (open/taken/closed) \\
created\_at & TIMESTAMP & Ngày tạo \\
\hline
\end{tabular}
\end{table}

\begin{table}[H]
\centering
\caption{Bảng Messages - Tin nhắn}
\begin{tabular}{|l|l|l|}
\hline
\textbf{Tên trường} & \textbf{Kiểu dữ liệu} & \textbf{Mô tả} \\
\hline
id & INT (PK) & Mã tin nhắn \\
from\_user & INT (FK) & Người gửi \\
to\_user & INT (FK) & Người nhận \\
post\_id & INT (FK) & Bài đăng liên quan \\
message & TEXT & Nội dung tin nhắn \\
created\_at & TIMESTAMP & Thời gian gửi \\
\hline
\end{tabular}
\end{table}

\begin{table}[H]
\centering
\caption{Bảng Verifications - Yêu cầu xác thực}
\begin{tabular}{|l|l|l|}
\hline
\textbf{Tên trường} & \textbf{Kiểu dữ liệu} & \textbf{Mô tả} \\
\hline
id & INT (PK) & Mã yêu cầu \\
user\_id & INT (FK) & Mã người dùng \\
full\_name & VARCHAR(150) & Họ tên đầy đủ \\
student\_code & VARCHAR(100) & Mã số sinh viên \\
class\_name & VARCHAR(100) & Lớp \\
document\_card & VARCHAR(255) & Ảnh thẻ sinh viên \\
status & ENUM & Trạng thái xử lý \\
admin\_note & TEXT & Ghi chú của admin \\
created\_at & TIMESTAMP & Ngày tạo \\
\hline
\end{tabular}
\end{table}

\begin{table}[H]
\centering
\caption{Bảng Reports - Báo cáo vi phạm}
\begin{tabular}{|l|l|l|}
\hline
\textbf{Tên trường} & \textbf{Kiểu dữ liệu} & \textbf{Mô tả} \\
\hline
id & INT (PK) & Mã báo cáo \\
reporter\_id & INT (FK) & Người báo cáo \\
target\_type & ENUM & Loại (post/comment) \\
target\_id & INT & ID đối tượng bị báo cáo \\
reason\_code & VARCHAR(50) & Mã lý do \\
message & TEXT & Nội dung chi tiết \\
status & ENUM & Trạng thái xử lý \\
created\_at & TIMESTAMP & Ngày tạo \\
\hline
\end{tabular}
\end{table}

\begin{table}[H]
\centering
\caption{Bảng Ratings - Đánh giá chất lượng dịch vụ}
\begin{tabular}{|l|l|l|}
\hline
\textbf{Tên trường} & \textbf{Kiểu dữ liệu} & \textbf{Mô tả} \\
\hline
id & INT (PK) & Khóa chính, tự tăng \\
user\_id & INT (FK) & ID người thực hiện đánh giá \\
rated\_user\_id & INT (FK) & ID người được đánh giá \\
rating & TINYINT & Điểm đánh giá (từ 1 đến 5 sao) \\
comment & TEXT & Nhận xét chi tiết về chuyên môn/thái độ \\
created\_at & TIMESTAMP & Thời gian gửi đánh giá \\
\hline
\end{tabular}
\end{table}

\begin{table}[H]
\centering
\caption{Bảng Appointments - Quản lý lịch hẹn y tế}
\begin{tabular}{|l|l|l|}
\hline
\textbf{Tên trường} & \textbf{Kiểu dữ liệu} & \textbf{Mô tả} \\
\hline
id & INT (PK) & Khóa chính, tự tăng \\
patient\_id & INT (FK) & ID người bệnh cần chăm sóc \\
student\_id & INT (FK) & ID sinh viên Y nhận việc \\
appointment\_date & DATETIME & Thời gian thực hiện chăm sóc \\
status & ENUM & Trạng thái lịch hẹn (pending/confirmed/completed/cancelled) \\
notes & TEXT & Ghi chú/yêu cầu lâm sàng đặc biệt \\
created\_at & TIMESTAMP & Thời gian tạo lịch hẹn \\
\hline
\end{tabular}
\end{table}

\begin{table}[H]
\centering
\caption{Bảng Notifications - Thông báo hệ thống}
\begin{tabular}{|l|l|l|}
\hline
\textbf{Tên trường} & \textbf{Kiểu dữ liệu} & \textbf{Mô tả} \\
\hline
id & INT (PK) & Khóa chính, tự tăng \\
user\_id & INT (FK) & ID người nhận thông báo \\
type & VARCHAR(50) & Phân loại thông báo (message/appointment/comment/verify) \\
title & VARCHAR(255) & Tiêu đề thông báo \\
message & TEXT & Nội dung thông báo chi tiết \\
link & VARCHAR(255) & Liên kết chuyển hướng nhanh khi click \\
is\_read & TINYINT & Trạng thái đã đọc (0: Chưa đọc, 1: Đã đọc) \\
created\_at & TIMESTAMP & Thời gian gửi thông báo \\
\hline
\end{tabular}
\end{table}

\begin{table}[H]
\centering
\caption{Bảng AI Conversations - Cuộc trò chuyện với Trợ lý ảo}
\begin{tabular}{|l|l|l|}
\hline
\textbf{Tên trường} & \textbf{Kiểu dữ liệu} & \textbf{Mô tả} \\
\hline
id & INT (PK) & Khóa chính, tự tăng \\
user\_id & INT (FK) & ID người dùng tương tác với AI \\
title & VARCHAR(255) & Tiêu đề phiên trò chuyện \\
created\_at & TIMESTAMP & Thời gian bắt đầu \\
\hline
\end{tabular}
\end{table}

\begin{table}[H]
\centering
\caption{Bảng AI Messages - Chi tiết tin nhắn hội thoại AI}
\begin{tabular}{|l|l|l|}
\hline
\textbf{Tên trường} & \textbf{Kiểu dữ liệu} & \textbf{Mô tả} \\
\hline
id & INT (PK) & Khóa chính, tự tăng \\
conversation\_id & INT (FK) & ID phiên trò chuyện tương ứng \\
role & ENUM & Vai trò người gửi (user / model) \\
content & TEXT & Nội dung câu hỏi hoặc tư vấn \\
source & VARCHAR(50) & Nguồn xử lý AI (gemini/huggingface) \\
created\_at & TIMESTAMP & Thời gian gửi tin nhắn \\
\hline
\end{tabular}
\end{table}

\begin{table}[H]
\centering
\caption{Bảng Knowledge Posts - Bài viết kiến thức y học}
\begin{tabular}{|l|l|l|}
\hline
\textbf{Tên trường} & \textbf{Kiểu dữ liệu} & \textbf{Mô tả} \\
\hline
id & INT (PK) & Khóa chính, tự tăng \\
user\_id & INT (FK) & ID tác giả bài viết \\
title & VARCHAR(255) & Tiêu đề bài viết chia sẻ sức khỏe \\
content & TEXT & Nội dung chuyên môn chi tiết \\
category & VARCHAR(100) & Chuyên khoa/Danh mục y học \\
tags & TEXT & Nhãn từ khóa tìm kiếm nhanh \\
is\_featured & TINYINT & Đánh dấu bài viết nổi bật trang chủ \\
view\_count & INT & Số lượt xem bài viết \\
created\_at & TIMESTAMP & Thời gian đăng bài \\
\hline
\end{tabular}
\end{table}

\begin{table}[H]
\centering
\caption{Bảng Activity Logs - Nhật ký hoạt động hệ thống}
\begin{tabular}{|l|l|l|}
\hline
\textbf{Tên trường} & \textbf{Kiểu dữ liệu} & \textbf{Mô tả} \\
\hline
id & INT (PK) & Khóa chính, tự tăng \\
user\_id & INT (FK) & ID người thực hiện hành động \\
action & VARCHAR(100) & Tên hành động thực hiện \\
description & TEXT & Mô tả chi tiết hành động \\
ip\_address & VARCHAR(45) & Địa chỉ IP thực hiện thao tác \\
user\_agent & TEXT & Thông tin trình duyệt/thiết bị \\
created\_at & TIMESTAMP & Thời gian ghi nhận nhật ký \\
\hline
\end{tabular}
\end{table}

\begin{table}[H]
\centering
\caption{Bảng Friendships - Quan hệ bạn bè giữa người dùng}
\begin{tabular}{|l|l|l|}
\hline
\textbf{Tên trường} & \textbf{Kiểu dữ liệu} & \textbf{Mô tả} \\
\hline
id & INT (PK) & Khóa chính, tự tăng \\
user\_id & INT (FK) & ID người gửi lời mời kết bạn \\
friend\_id & INT (FK) & ID người nhận lời mời kết bạn \\
status & ENUM & Trạng thái kết bạn (pending/accepted/blocked) \\
created\_at & TIMESTAMP & Thời gian gửi lời mời \\
accepted\_at & TIMESTAMP & Thời gian chấp nhận kết bạn \\
\hline
\end{tabular}
\end{table}

\section{Cài đặt hệ thống}

\subsection{Cấu trúc thư mục dự án}

\begin{verbatim}
DACN1/
├── assets/
│   ├── css/
│   │   └── style.css
│   └── js/
├── uploads/
│   ├── student_cards/
│   └── verification_docs/
├── config.php
├── header.php
├── footer.php
├── index.php
├── login.php
├── register.php
├── dashboard_patient.php
├── dashboard_student.php
├── create_recruitment.php
├── create_application.php
├── view_post.php
├── chat.php
├── conversations.php
├── admin_dashboard.php
├── admin_users.php
├── admin_posts.php
├── admin_verifications.php
├── admin_reports.php
└── database.sql
\end{verbatim}

\subsection{Cấu hình kết nối cơ sở dữ liệu}

File \texttt{config.php} chứa thông tin kết nối database và các hàm tiện ích:

\begin{lstlisting}[language=PHP]
<?php
define('DB_HOST', '127.0.0.1');
define('DB_NAME', 'dacn1_db');
define('DB_USER', 'root');
define('DB_PASS', '');

$pdo = new PDO(
    'mysql:host='.DB_HOST.';dbname='.DB_NAME,
    DB_USER, DB_PASS,
    [PDO::ATTR_ERRMODE => PDO::ERRMODE_EXCEPTION]
);
\end{lstlisting}

\subsection{Xử lý đăng ký và phân quyền}

Hệ thống tự động phân quyền dựa trên email:
\begin{itemize}
    \item Email kết thúc bằng \texttt{edu.vn} $\rightarrow$ Vai trò: Sinh viên
    \item Email khác $\rightarrow$ Vai trò: Bệnh nhân
\end{itemize}

\subsection{Xác thực sinh viên}

Quy trình xác thực sinh viên:
\begin{enumerate}
    \item Sinh viên gửi yêu cầu xác thực kèm ảnh thẻ sinh viên
    \item Admin xem xét và duyệt/từ chối yêu cầu
    \item Nếu được duyệt, sinh viên có thể đăng tin ứng tuyển
\end{enumerate}

\subsection{Hệ thống tin nhắn}

Hệ thống hỗ trợ nhắn tin trực tiếp giữa người dùng:
\begin{itemize}
    \item Tạo cuộc hội thoại mới khi nhắn tin lần đầu
    \item Lưu trữ lịch sử tin nhắn
    \item Hiển thị trạng thái online của người dùng
\end{itemize}


% ==================== CHƯƠNG 4: KẾT QUẢ NGHIÊN CỨU ====================
\chapter{KẾT QUẢ NGHIÊN CỨU}

\section{Giao diện người dùng}

\subsection{Trang chủ}

Trang chủ hiển thị danh sách các tin đăng mới nhất, cho phép người dùng tìm kiếm và lọc theo danh mục, khu vực.

\begin{figure}[H]
    \centering
    \includegraphics[width=0.9\textwidth]{trangchu}
    \caption{Giao diện trang chủ}
    \label{fig:trangchu}
\end{figure}

\subsection{Đăng ký và Đăng nhập}

\begin{figure}[H]
    \centering
    \includegraphics[width=0.8\textwidth]{dangky}
    \caption{Giao diện đăng ký tài khoản}
    \label{fig:dangky}
\end{figure}

\begin{figure}[H]
    \centering
    \includegraphics[width=0.8\textwidth]{dangnhap}
    \caption{Giao diện đăng nhập}
    \label{fig:dangnhap}
\end{figure}

\subsection{Trang Bệnh nhân}

\subsubsection{Bảng điều khiển}

\begin{figure}[H]
    \centering
    \includegraphics[width=0.9\textwidth]{dashboard_benhnhan}
    \caption{Bảng điều khiển trang bệnh nhân}
    \label{fig:dashboard_patient}
\end{figure}

\begin{figure}[H]
    \centering
    \includegraphics[width=0.3\textwidth]{menu_benhnhan}
    \caption{Menu trang bệnh nhân}
    \label{fig:menu_patient}
\end{figure}

\subsubsection{Tạo tin tuyển dụng}

\begin{figure}[H]
    \centering
    \includegraphics[width=0.9\textwidth]{taotintuyen}
    \caption{Giao diện tạo tin tuyển dụng}
    \label{fig:create_recruitment}
\end{figure}

\subsubsection{Đánh giá sinh viên}

\begin{figure}[H]
    \centering
    \includegraphics[width=0.8\textwidth]{danhgia}
    \caption{Giao diện đánh giá sinh viên}
    \label{fig:rating}
\end{figure}

\subsubsection{Yêu thích và Lịch sử}

\begin{figure}[H]
    \centering
    \includegraphics[width=0.9\textwidth]{yeuthich_benhnhan}
    \caption{Giao diện tin yêu thích - Bệnh nhân}
    \label{fig:favorite_patient}
\end{figure}

\begin{figure}[H]
    \centering
    \includegraphics[width=0.9\textwidth]{lichsu_benhnhan}
    \caption{Giao diện lịch sử nhận việc - Bệnh nhân}
    \label{fig:history_patient}
\end{figure}

\subsection{Trang Sinh viên}

\subsubsection{Bảng điều khiển}

\begin{figure}[H]
    \centering
    \includegraphics[width=0.9\textwidth]{dashboard_sinhvien}
    \caption{Bảng điều khiển trang sinh viên}
    \label{fig:dashboard_student}
\end{figure}

\begin{figure}[H]
    \centering
    \includegraphics[width=0.3\textwidth]{menu_sinhvien}
    \caption{Menu trang sinh viên}
    \label{fig:menu_student}
\end{figure}

\subsubsection{Cập nhật hồ sơ}

\begin{figure}[H]
    \centering
    \includegraphics[width=0.8\textwidth]{capnhathoso}
    \caption{Form cập nhật hồ sơ sinh viên}
    \label{fig:update_profile}
\end{figure}

\subsubsection{Xin xác thực và cấp quyền đăng tin}

\begin{figure}[H]
    \centering
    \includegraphics[width=0.8\textwidth]{xinxacthuc}
    \caption{Form xin xác thực tài khoản}
    \label{fig:request_verify}
\end{figure}

\begin{figure}[H]
    \centering
    \includegraphics[width=0.8\textwidth]{xincapquyen}
    \caption{Form xin cấp quyền đăng tin}
    \label{fig:request_post}
\end{figure}

\begin{figure}[H]
    \centering
    \includegraphics[width=0.8\textwidth]{thongbao_capquyen}
    \caption{Thông báo xin cấp quyền đăng tin}
    \label{fig:notify_post}
\end{figure}

\subsubsection{Tạo tin ứng tuyển}

\begin{figure}[H]
    \centering
    \includegraphics[width=0.9\textwidth]{taotinungtuyen}
    \caption{Giao diện tạo tin ứng tuyển}
    \label{fig:create_application}
\end{figure}

\subsubsection{Lịch sử nhận việc}

\begin{figure}[H]
    \centering
    \includegraphics[width=0.9\textwidth]{lichsu_sinhvien}
    \caption{Lịch sử nhận việc - Sinh viên}
    \label{fig:history_student}
\end{figure}

\subsection{Chức năng chung}

\subsubsection{Hội thoại và Tin nhắn}

\begin{figure}[H]
    \centering
    \includegraphics[width=0.9\textwidth]{hoithoai}
    \caption{Giao diện hội thoại}
    \label{fig:chat}
\end{figure}

\subsubsection{Bạn bè}

\begin{figure}[H]
    \centering
    \includegraphics[width=0.9\textwidth]{banbe}
    \caption{Giao diện trang bạn bè}
    \label{fig:friends}
\end{figure}

\subsubsection{Yêu thích}

\begin{figure}[H]
    \centering
    \includegraphics[width=0.9\textwidth]{yeuthich}
    \caption{Giao diện bài tuyển yêu thích}
    \label{fig:favorites}
\end{figure}

\subsubsection{Thông báo}

\begin{figure}[H]
    \centering
    \includegraphics[width=0.8\textwidth]{thongbao}
    \caption{Giao diện thông báo}
    \label{fig:notification}
\end{figure}

\subsubsection{Hỗ trợ tài khoản}

\begin{figure}[H]
    \centering
    \includegraphics[width=0.9\textwidth]{hotro}
    \caption{Giao diện hỗ trợ tài khoản}
    \label{fig:support}
\end{figure}

\subsection{Trang Quản trị}

\subsubsection{Bảng điều khiển Admin}

\begin{figure}[H]
    \centering
    \includegraphics[width=0.9\textwidth]{admin_dashboard}
    \caption{Bảng điều khiển trang quản trị}
    \label{fig:admin_dashboard}
\end{figure}

\begin{figure}[H]
    \centering
    \includegraphics[width=0.3\textwidth]{menu_admin}
    \caption{Menu trang quản trị}
    \label{fig:menu_admin}
\end{figure}

\subsubsection{Quản lý người dùng}

\begin{figure}[H]
    \centering
    \includegraphics[width=0.9\textwidth]{admin_users}
    \caption{Giao diện quản lý người dùng}
    \label{fig:admin_users}
\end{figure}

\subsubsection{Quản lý bài viết}

\begin{figure}[H]
    \centering
    \includegraphics[width=0.9\textwidth]{admin_posts}
    \caption{Giao diện quản lý bài viết}
    \label{fig:admin_posts}
\end{figure}

\subsubsection{Xác minh sinh viên}

\begin{figure}[H]
    \centering
    \includegraphics[width=0.9\textwidth]{admin_xacminh}
    \caption{Giao diện xác minh sinh viên}
    \label{fig:admin_verify}
\end{figure}

\subsubsection{Quản lý tin đăng}

\begin{figure}[H]
    \centering
    \includegraphics[width=0.9\textwidth]{admin_tintuyen}
    \caption{Quản lý tin tuyển dụng}
    \label{fig:admin_recruitment}
\end{figure}

\begin{figure}[H]
    \centering
    \includegraphics[width=0.9\textwidth]{admin_tinungtuyen}
    \caption{Quản lý tin ứng tuyển}
    \label{fig:admin_application}
\end{figure}

\subsubsection{Xem khiếu nại}

\begin{figure}[H]
    \centering
    \includegraphics[width=0.9\textwidth]{admin_khieunai}
    \caption{Giao diện xem khiếu nại}
    \label{fig:admin_reports}
\end{figure}

\subsubsection{Gửi thông báo}

\begin{figure}[H]
    \centering
    \includegraphics[width=0.8\textwidth]{admin_guithongbao}
    \caption{Giao diện gửi thông báo}
    \label{fig:admin_notify}
\end{figure}

\subsubsection{Yêu thích - Admin}

\begin{figure}[H]
    \centering
    \includegraphics[width=0.9\textwidth]{admin_yeuthich}
    \caption{Giao diện yêu thích - Quản trị}
    \label{fig:admin_favorites}
\end{figure}

\subsubsection{Lịch sử nhận việc - Admin}

\begin{figure}[H]
    \centering
    \includegraphics[width=0.9\textwidth]{admin_lichsu}
    \caption{Giao diện lịch sử nhận việc - Quản trị}
    \label{fig:admin_history}
\end{figure}


\section{Đánh giá kết quả}

\subsection{Các chức năng đã hoàn thành}

\begin{table}[H]
\centering
\caption{Bảng đánh giá các chức năng}
\begin{tabular}{|c|l|c|}
\hline
\textbf{STT} & \textbf{Chức năng} & \textbf{Trạng thái} \\
\hline
1 & Đăng ký/Đăng nhập/Đăng xuất & Hoàn thành \\
2 & Phân loại người dùng tự động & Hoàn thành \\
3 & Xác minh sinh viên Y & Hoàn thành \\
4 & Đăng tin tuyển dụng (bệnh nhân) & Hoàn thành \\
5 & Đăng tin ứng tuyển (sinh viên) & Hoàn thành \\
6 & Tìm kiếm và lọc tin & Hoàn thành \\
7 & Nhắn tin trực tiếp & Hoàn thành \\
8 & Kết bạn & Hoàn thành \\
9 & Đánh giá người dùng & Hoàn thành \\
10 & Lưu tin yêu thích & Hoàn thành \\
11 & Báo cáo vi phạm & Hoàn thành \\
12 & Quản lý người dùng (admin) & Hoàn thành \\
13 & Kiểm duyệt bài viết (admin) & Hoàn thành \\
14 & Xử lý khiếu nại (admin) & Hoàn thành \\
15 & Gửi thông báo (admin) & Hoàn thành \\
16 & Gọi Video trực tiếp thời gian thực (WebRTC) & Hoàn thành \\
17 & Trợ lý ảo AI tư vấn y học (Hugging Face) & Hoàn thành \\
18 & Hệ thống đặt và quản lý lịch hẹn (Appointments) & Hoàn thành \\
19 & Cổng kiến thức sức khỏe cộng đồng (Knowledge Posts) & Hoàn thành \\
20 & Nhật ký hoạt động an ninh hệ thống (Activity Logs) & Hoàn thành \\
\hline
\end{tabular}
\end{table}

\subsection{Ưu điểm}
\begin{itemize}
    \item \textbf{Giao diện thân thiện:} Thiết đơn giản, trực quan, dễ sử dụng cho người bệnh và người cao tuổi.
    \item \textbf{Responsive:} Hiển thị và vận hành tối ưu trên cả desktop lẫn các thiết bị di động.
    \item \textbf{Độ tin cậy cao:} Cơ chế xác thực thẻ sinh viên nghiêm ngặt được phê duyệt bởi Admin giúp bảo vệ cộng đồng.
    \item \textbf{Giao tiếp đa chiều:} Cuộc gọi Video Call qua WebRTC giúp hai bên đối chiếu trực quan thông tin và kiểm tra lâm sàng vết thương từ xa.
    \item \textbf{Tương tác thông minh:} Tích hợp Trợ lý ảo AI Chatbot (BlenderBot/Gemini) hỗ trợ người dùng giải đáp y khoa cơ bản 24/7.
    \item \textbf{Quản lý chuyên nghiệp:} Tự động hóa quy trình tạo lịch hẹn chăm sóc và nhật ký kiểm toán hệ thống.
\end{itemize}

\subsection{Hạn chế}
\begin{itemize}
    \item Chưa phát triển ứng dụng di động dạng Native (iOS/Android).
    \item Chưa tích hợp cổng thanh toán trực tuyến tự động (MoMo, VNPay...).
    \item Chưa tích hợp định vị GPS bản đồ (Google Maps) để tìm kiếm theo khoảng cách địa lý.
\end{itemize}

% ==================== CHƯƠNG 5: KẾT LUẬN VÀ HƯỚNG PHÁT TRIỂN ====================
\chapter{KẾT LUẬN VÀ HƯỚNG PHÁT TRIỂN}

\section{Kết luận}

Sau quá trình nghiên cứu và thực hiện, đồ án "Xây dựng Website kết nối y tế giữa bệnh nhân và sinh viên Y khoa" đã đạt được các kết quả sau:

\subsection{Về mặt lý thuyết}
\begin{itemize}
    \item Nắm vững kiến thức về phát triển web với PHP và MySQL
    \item Hiểu rõ quy trình phân tích, thiết kế và xây dựng hệ thống thông tin
    \item Áp dụng được các mô hình UML trong thiết kế phần mềm
    \item Hiểu về bảo mật web cơ bản
\end{itemize}

\subsection{Về mặt thực tiễn}
\begin{itemize}
    \item Xây dựng thành công website với đầy đủ chức năng cơ bản kết nối bệnh nhân và sinh viên Y khoa.
    \item Hệ thống phân quyền hoạt động chính xác và bảo mật.
    \item Giao diện thân thiện, tương thích tốt trên cả máy tính và thiết bị di động (Responsive).
    \item Tích hợp thành công cơ chế gọi điện thoại và cuộc gọi Video trực tuyến thời gian thực bằng công nghệ WebRTC.
    \item Tích hợp Trợ lý ảo AI Chatbot tư vấn y học cơ bản 24/7 (sử dụng APIs Hugging Face/Gemini).
    \item Hoàn thành hệ thống đặt lịch hẹn chăm sóc sức khỏe tự động (Appointments) và Nhật ký hoạt động (Activity Logs).
    \item Hệ thống xác thực sinh viên qua thẻ sinh viên hoạt động tin cậy dưới sự kiểm duyệt của Admin.
    \item Các chức năng tìm kiếm, nhắn tin, kết bạn, đánh giá hoạt động ổn định.
\end{itemize}

\subsection{Những đóng góp của đồ án}
\begin{itemize}
    \item Tạo ra nền tảng kết nối trực tuyến an toàn và đáng tin cậy giữa bệnh nhân và sinh viên Y khoa.
    \item Giúp bệnh nhân dễ dàng tiếp cận dịch vụ hỗ trợ y tế chất lượng tại nhà với chi phí hợp lý.
    \item Tạo cơ hội thực hành lâm sàng thực tế ngoài bệnh viện và tạo thêm thu nhập chính đáng cho sinh viên Y khoa.
    \item Góp phần giảm tải áp lực cho hệ thống y tế công cộng và hỗ trợ y tế cộng đồng.
\end{itemize}

\section{Hướng phát triển}

Để hoàn thiện và mở rộng hệ thống, các hướng phát triển trong tương lai bao gồm:

\subsection{Về chức năng}
\begin{itemize}
    \item Phát triển ứng dụng mobile native (Android/iOS) bằng React Native hoặc Flutter.
    \item Tích hợp cổng thanh toán trực tuyến tự động (MoMo, VNPay...) và cơ chế đặt cọc giữ tiền.
    \item Áp dụng thuật toán Machine Learning để xây dựng hệ thống gợi ý sinh viên Y khoa phù hợp nhất dựa trên bệnh án của người bệnh.
    \item Tích hợp giao thức WebSocket để hỗ trợ tin nhắn và thông báo thời gian thực (Real-time).
    \item Tích hợp Google Maps API để hỗ trợ định vị vị trí và tìm kiếm xung quanh.
\end{itemize}

\subsection{Về kỹ thuật}
\begin{itemize}
    \item Triển khai ứng dụng trên các hệ thống Cloud Server thực tế (AWS, Google Cloud).
    \item Tối ưu hóa hiệu năng, bảo mật thông tin y tế người dùng theo chuẩn HIPAA.
    \item Áp dụng framework PHP (như Laravel) để chuẩn hóa mã nguồn và dễ bảo trì.
\end{itemize}

\subsection{Về nghiệp vụ}
\begin{itemize}
    \item Mở rộng đối tượng người dùng đến điều dưỡng chuyên nghiệp, bác sĩ gia đình.
    \item Hợp tác liên kết trực tiếp với các trường đại học y dược và các bệnh viện để mở rộng quy mô.
\end{itemize}


% ==================== TÀI LIỆU THAM KHẢO ====================
\newpage
\addcontentsline{toc}{chapter}{DANH MỤC TÀI LIỆU THAM KHẢO}
\begin{center}
\textbf{\Large DANH MỤC TÀI LIỆU THAM KHẢO}
\end{center}
\vspace{1cm}

\begin{enumerate}[label={[\arabic*]}]
    \item PHP Documentation. (2024). \textit{PHP Manual}. https://www.php.net/manual/en/
    
    \item MySQL Documentation. (2024). \textit{MySQL 8.0 Reference Manual}. https://dev.mysql.com/doc/refman/8.0/en/
    
    \item Bootstrap Team. (2024). \textit{Bootstrap Documentation}. https://getbootstrap.com/docs/
    
    \item W3Schools. (2024). \textit{HTML Tutorial}. https://www.w3schools.com/html/
    
    \item W3Schools. (2024). \textit{CSS Tutorial}. https://www.w3schools.com/css/
    
    \item Mozilla Developer Network. (2024). \textit{JavaScript Guide}. https://developer.mozilla.org/en-US/docs/Web/JavaScript/Guide
    
    \item Apache Friends. (2024). \textit{XAMPP Documentation}. https://www.apachefriends.org/
    
    \item Nguyễn Văn A. (2023). \textit{Giáo trình Phát triển ứng dụng Web}. NXB Đại học Quốc gia.
    
    \item Trần Văn B. (2022). \textit{Phân tích và thiết kế hệ thống thông tin}. NXB Giáo dục.
    
    \item Object Management Group. (2024). \textit{UML Specification}. https://www.omg.org/spec/UML/
\end{enumerate}

% ==================== PHỤ LỤC ====================
\newpage
\addcontentsline{toc}{chapter}{PHỤ LỤC}
\begin{center}
\textbf{\Large PHỤ LỤC}
\end{center}
\vspace{1cm}

\section*{Phụ lục A: Hướng dẫn cài đặt}

\subsection*{Yêu cầu hệ thống}
\begin{itemize}
    \item XAMPP (Apache + PHP + MySQL) phiên bản 7.4 trở lên
    \item Trình duyệt web hiện đại (Chrome, Firefox, Edge)
    \item Hệ điều hành: Windows 10/11
\end{itemize}

\subsection*{Các bước cài đặt}
\begin{enumerate}
    \item Cài đặt XAMPP từ https://www.apachefriends.org/
    \item Copy thư mục dự án vào \texttt{C:\textbackslash xampp\textbackslash htdocs\textbackslash DACN1}
    \item Khởi động Apache và MySQL trong XAMPP Control Panel
    \item Mở phpMyAdmin: http://localhost/phpmyadmin
    \item Tạo database mới với tên \texttt{dacn1\_db}
    \item Import file \texttt{database.sql}
    \item Cấu hình file \texttt{config.php} nếu cần
    \item Truy cập http://localhost/DACN1
\end{enumerate}

\subsection*{Tài khoản mặc định}
\begin{itemize}
    \item Admin: Tạo qua file \texttt{create\_admin.php}
    \item Người dùng: Đăng ký qua trang \texttt{register.php}
\end{itemize}

\section*{Phụ lục B: Danh sách các file chính}

\begin{table}[H]
\centering
\begin{tabular}{|l|l|}
\hline
\textbf{File} & \textbf{Mô tả} \\
\hline
config.php & Cấu hình database và hàm tiện ích \\
header.php & Header chung cho các trang \\
footer.php & Footer chung cho các trang \\
index.php & Trang chủ, danh sách tin \\
login.php & Đăng nhập \\
register.php & Đăng ký \\
logout.php & Đăng xuất \\
dashboard\_patient.php & Trang bệnh nhân \\
dashboard\_student.php & Trang sinh viên \\
create\_recruitment.php & Tạo tin tuyển dụng \\
create\_application.php & Tạo tin ứng tuyển \\
view\_post.php & Xem chi tiết tin \\
chat.php & Nhắn tin \\
conversations.php & Danh sách hội thoại \\
friends.php & Quản lý bạn bè \\
favorites.php & Tin yêu thích \\
profile.php & Hồ sơ cá nhân \\
edit\_profile.php & Chỉnh sửa hồ sơ \\
request\_verification.php & Xin xác thực thẻ sinh viên \\
request\_posting\_permission.php & Xin quyền đăng tin \\
admin\_dashboard.php & Trang quản trị \\
admin\_users.php & Quản lý người dùng \\
admin\_posts.php & Quản lý bài viết \\
admin\_verifications.php & Xử lý xác thực \\
admin\_reports.php & Xem báo cáo \\
admin\_send\_notification.php & Gửi thông báo \\
video\_call.php & Cuộc gọi video trực tuyến thời gian thực (WebRTC) \\
call\_signaling.php & Xử lý báo hiệu tín hiệu cuộc gọi video \\
ai\_assistant.php & Trang Trợ lý ảo AI tư vấn \\
chatbot\_widget.php & Widget hiển thị bong bóng chat AI \\
create\_appointment.php & Tạo lịch hẹn chăm sóc tự động \\
appointment\_details.php & Chi tiết lịch hẹn y tế \\
assignment\_history.php & Lịch sử phân công lịch hẹn \\
qr\_scanner\_widget.php & Widget quét mã QR xác thực lịch hẹn \\
\hline
\end{tabular}
\end{table}

\section*{Phụ lục C: Hướng dẫn đổi tên file ảnh cho Overleaf}

Để sử dụng ảnh trong Overleaf, bạn cần đổi tên các file ảnh từ thư mục "Ảnh Đồ Án" sang tên không dấu và upload vào thư mục \texttt{images/} trên Overleaf:

\begin{table}[H]
\centering
\small
\begin{tabular}{|l|l|}
\hline
\textbf{Tên file gốc} & \textbf{Tên file mới} \\
\hline
Giao diện trang chủ.png & trangchu.png \\
Giao diện đăng kí.png & dangky.png \\
Giao diện đăng nhập.png & dangnhap.png \\
Giao diện bảng điều khiển trang bệnh nhân.png & dashboard\_benhnhan.png \\
menu trang bệnh nhân.png & menu\_benhnhan.png \\
Giao diện tạo tin tuyển .png & taotintuyen.png \\
Giao diện đánh giá trang bệnh nhân.png & danhgia.png \\
Giao diện yêu thích trang bệnh nhân.png & yeuthich\_benhnhan.png \\
Giao diện lịch sử nhận việc trang bệnh nhân.png & lichsu\_benhnhan.png \\
Giao diện bảng điều khiển trang sinh viên.png & dashboard\_sinhvien.png \\
Menu trang sinh viên.png & menu\_sinhvien.png \\
Giao diện form cập nhật hồ sơ trang sinh viên.png & capnhathoso.png \\
Giao diện form xin xác thực tài khoản.png & xinxacthuc.png \\
Giao diện form xin cấp quyền đăng tin.png & xincapquyen.png \\
Giao diện thông báo xin cấp quyền đăng tin.png & thongbao\_capquyen.png \\
Giao diện tạo tin ứng tuyển.png & taotinungtuyen.png \\
Giao diện lịch sử nhận việc trang sinh viên.png & lichsu\_sinhvien.png \\
Giao diện hội thoại.png & hoithoai.png \\
Giao diện trang bạn bè.png & banbe.png \\
Giao diện trang bài tuyển yêu thích.png & yeuthich.png \\
Giao diện thông báo.png & thongbao.png \\
Giao diện trang hỗ trợ tài khoản.png & hotro.png \\
Giao diện bảng điều khiển trang quản trị.png & admin\_dashboard.png \\
Menu trang quản trị.png & menu\_admin.png \\
Giao diện quản lý người dùng.png & admin\_users.png \\
Giao diện quản lý bài viết.png & admin\_posts.png \\
Giao diện xác minh sinh viên.png & admin\_xacminh.png \\
Giao diện đăng tin tuyển trang quản trị.png & admin\_tintuyen.png \\
Giao diện đăng tin ứng tuyển trang quản trị.png & admin\_tinungtuyen.png \\
Giao diện xem khiếu nại.png & admin\_khieunai.png \\
Giao dện gửi thông báo .png & admin\_guithongbao.png \\
Giao diện yêu thích trang quản trị.png & admin\_yeuthich.png \\
Giao diện trang lịch sử nhận việc.png & admin\_lichsu.png \\
\hline
\end{tabular}
\end{table}

\textbf{Lưu ý:} Sau khi đổi tên, upload tất cả ảnh vào thư mục \texttt{images/} trên Overleaf. File usecase.SVG và erd.SVG cũng cần được upload vào thư mục \texttt{images/} với tên \texttt{usecase.png} và \texttt{erd.png} (chuyển đổi sang PNG nếu cần).

\end{document}
